\documentclass[12pt, oneside]{article}   	% use "amsart" instead of "article" for AMSLaTeX format
\usepackage{geometry}                		% See geometry.pdf to learn the layout options. There are lots.
\geometry{letterpaper}                   		% ... or a4paper or a5paper or ... 
%\geometry{landscape}                		% Activate for for rotated page geometry
%\usepackage[parfill]{parskip}    		% Activate to begin paragraphs with an empty line rather than an indent
\usepackage{graphicx}				% Use pdf, png, jpg, or eps� with pdflatex; use eps in DVI mode
								% TeX will automatically convert eps --> pdf in pdflatex		
\usepackage{amssymb}


\newcommand{\bfiv}{{\bf{{[5 points]}}~~}}
%\newcommand{\bfiv}{{\bf{\ul{[5 points]}}~}}

\input{/Users/colmocinneide/Desktop/Latex/latex_defs}

 \setlength{\topmargin}{-.75in}
\setlength{\textheight}{9.5in}
\setlength{\oddsidemargin}{-.5in}
\setlength{\textwidth}{7in}
\pagestyle{myheadings} 
\markright{ \bf \Large \em  NAME:  {\hspace{3in}  }ID: {\hspace{1.4 in} }}

\begin{document}

\centerline{\bf Quiz 1(Wednesday, March 22, 2017) }
\vspace{.1in}

\noindent
{\em Please write your name and student ID at the top of each page.
Please circle your answers or write them in the boxes provided.} 
  Worth 10\% of your total grade for the course.  You have 30 minutes to compete the quiz. 
  Electronic devices such as phones, calculators and laptops, are not permitted.  
  Finish by 7:20 and return completed quiz paper  to the TAs. 
 
 \xbe


 
 
 
 
% \item In the Bayesian approach, there are distributions of parameters and distributions of returns on assets.  Circle the distributions of returns in this list: \\\
% Statistical model, Prior distribution, posterior distribution, predictive distribution


 \item 
 Forty percent of a population has a certain disease.  A diagnostic test correctly 
 detects the disease 80\% of the time but produces 30\% false positives (i.e., incorrectly indicates that the person has the disease).  
 Circle the correct answers.
 \xbe
 \item \bfiv
 If a person is selected at random, what is the probability that his test result will be a {\em false positive}?
 {\em (i) 5\%
 (ii) 14\%
  \ul{\bf (iii) 18\%}
   (iv) 27\%
    (v) 32\%}
    \item \bfiv
         Given that the person gets a
 positive test result, what is the conditional probability he has the disease?
 {\em (i) 4\%
 (ii) 8\%
  (iii) 16\%
   (iv) 32\%
     \ul{\bf (iii)  (v) 64\%}}
    \xee
    
    
 \item \bfiv Continuing the previous question, classify the following probabilities as  
describing the prior distribution, the statistical model, or the posterior distribution.
The unknown parameter here is {\em whether or not the selected individual has the disease.}
 \xbe
 \item Forty percent of a population have a certain disease.  
 \item A diagnostic test correctly 
 detects the disease 80\% of the time.  
 \item The probability that  the selected person has the disease given that 
 he tests positive. 
 \xee
 Indicate your answers by writing a letter (a)--(c) in each box.\\
\framebox{   \parbox[t]{.5in}{  \rule{0in}{.25in} (c)}} Posterior distribution
\framebox{   \parbox[t]{.5in}{  \rule{0in}{.25in} (b)}} Statistical model 
 \framebox{   \parbox[t]{.5in}{  \rule{0in}{.25in} (a)}} Prior distribution
 

    
     \item Annual  returns of a certain investment strategy  follow a
     normal distribution with known volatility 12\%.  
     Our prior distribution of the expected excess return (EER) is normal with mean 2\% standard deviation 4\%.  An observation $X$ has expected value equal to the EER and has standard deviation 5\%.   When we observe $X$ we find its value to be $x=4\%$.
     \xbe
%     \item \bfiv
%     What is the posterior expected value of the EER?  This can be answered  
%      without calculation.
%      {\em (i) 1.54\% (ii) 2.23\% (iii) 3.39\% (iv) 4.19\%}
      \item \bfiv
     What is the posterior standard deviation of the EER? \\  
      {\em      \ul{\bf  (i) 3.12\% } (ii) 4.23\% (iii) 4.9\% (iv) 5.1\%}
      \item \bfiv
     What is the predictive volatility of future returns?  \\ %This can be answered   without calculation.
      {\em (i) 8.5\% (ii) 11.9\%      \ul{\bf (iii) 12.4\% } (iv) 20.4\%}
      \xee
      
      
      

     
          \item \bfiv An investor holds 40\% of his wealth in stocks and 60\% in riskless bonds.
          Assuming that (a) the excess return of stocks over riskless bonds is 3\%, 
          (b) the volatility of stocks is 15\%, and (c)  the investor holds the mean-variance optimal portfolio of stocks and riskless bonds, what is his risk aversion?\\
 {\em (i) 1.333
 (ii) 2.333
      \ul{\bf  (iii) 3.333}
   (iv) 4.333
    (v) 5.333}

     \item The benchmark for this question  is \ul {40\% stocks and 60\% risky bonds}.  
     Risk aversion is $\lambda = 4$.  The covariance matrix
     of stocks and bonds is
     $$
\left(\begin{array}{cc}
.01 & 0\\
0 & .0025 
\end{array}
\right).
$$
\xbe
\item \bfiv Find the correlation of stocks and bonds.
 {\em      \ul{\bf (i) 0}
 (ii) .2
  (iii) .3
   (iv) .4
    (v) .5}
\item \bfiv Find the volatility of stocks.   {\em      \ul{\bf  (i) 10\%}
 (ii) 15\%
  (iii) 20\%
   (iv) 25\%
    (v) 30\%}
\item \bfiv What is the (Black-Litterman) implied expected excess return of stocks?  \\   {\em (i) 0\%
 (ii) 0.8\%
  (iii) 1.0\%
   (iv) 1.3\%
       \ul{\bf   (v) 1.6\%} }

\item \bfiv What is the (Black-Litterman) implied expected excess return of bonds?   \\  {\em (i) 0.4\%
      \ul{\bf (ii) 0.6\%}
  (iii) .08\%
   (iv) 1.2\%
    (v) 2.4\%}

\item \bfiv Find the benchmark volatility. {\em (i) 3\%
     \ul{\bf  (ii) 5\%}
  (iii) 7\%
   (iv) 9.123\%
    (v) 11\%}
\item \bfiv What is the contribution to risk of stocks (that is, the contribution of stocks to the risk of the benchmark)?
{\em (i) 1.2\%
 (ii) 2.2\%
     \ul{\bf   (iii) 3.2\%}
   (iv) 4.2\%
    (v) 5.2\%}
\item \bfiv What is the contribution to  risk of bonds (that is, the contribution of bonds to the risk of the benchmark)?
{\em     \ul{\bf  (i) 1.8\%}
 (ii) 2.8\%
  (iii) 3.8\%
   (iv) 4.8\%
    (v) 5.8\%}
 \xee
 
      
      \item \bfiv In the following list, associate each asset class  with the typical annual volatility of that asset class.  Indicate your answers by writing a letter (a)--(d) in each box.\\
{\bf Asset class:} \hfill {\em (a) Currency (b) Equity (c) Long-term US bonds (d) Short-term US bonds.}\\
$~$  {\bf Risk: } \hspace{.5in} \framebox{   \parbox[t]{.5in}{  \rule{0in}{.25in} (d) }} 1\% 
\framebox{   \parbox[t]{.5in}{  \rule{0in}{.25in}(c) }} 5\% 
\framebox{   \parbox[t]{.5in}{  \rule{0in}{.25in} (a)}} 10\%
\framebox{   \parbox[t]{.5in}{  \rule{0in}{.25in} (b)}} 15\%

      \item {\bf[Extra credit]} I have regressed the returns $y$ of a mutual fund on a broad market index $x$ in order to do a factor risk   analysis.
%      explaining the risk of the mutual fund
%     in terms of {\em their exposure to broad equity market risk} and {\em idiosyncratic
%     risk (that is, risk uncorrelated with the broad market)}.  
     The regression coefficient is $\beta = .8$
      and the idiosyncratic volatility is .12.
     Assume that broad equity market risk is 20\%.
%     \xbe
%     \item What are the {\em factors} in this example.  Circle the appropriate terms: $x$, $y$,  idiosyncratic risk.
%     \xee
%     The following questions concern an attribution of the risk of the mutual fund to two sources: 
%     (1) broad equity market risk and (2) idiosyncratic risk.
     \xbe
     \item \bfiv What is the volatility of the mutual fund? {\em (a) 16\% (b) 18\%      \ul{\bf (c) 20\%} (d) 22\% (e) 24\%}
     \item \bfiv What is the contribution to risk of equity risk to the risk of the mutual fund, as a percentage of the risk of the mutual fund?  (This is known as PCTR or { \em percent contribution to risk.})
     {\em (a) 36\% (b) 50\%      \ul{\bf  (c) 64\%} (d) 75\% (e) 90\%}
     \xee

     
 \xee
 
 
 
 {\bf Formulas:}  Not all formulas are provided, nor is every formula provided necessarily needed.


\xbi

\item Bayes' rule: $P(A|B)
={ \displaystyle P(B|A)P(A) \over \displaystyle P(B|A)P(A)  + P(B|\bar A)P(\bar A)}$
\item Posterior  (or predictive) mean in the simple normal case: 
$\tilde m = (m/s^2+ x /  \sigma^2)/
(1/s^2+ 1/ \sigma^2)$, where $m$ is the prior mean, $s$ is the prior standard deviation, 
$x$ is the observation
and
$\sigma$ is the standard deviation of the observation.
\item The posterior variance in this case: $\tilde s^2 = (1/s^2+ 1/  \sigma^2)^{-1}$.
\item Predictive distribution in this case: normal$(\tilde m,\sigma_0^2+\tilde s^2)$,
where $\sigma_0$ is the volatility of return.
\item Factor risk model: $\Sigma = B \Sigma_X  B^T + \Omega$  ($B$: multiple regression coefficients of
returns on factors; $\Sigma_X$: covariance matrix of factors; $\Omega$: covariance of residuals.)
\item  Given the decomposition
$Y=\sum_{i=1}^n X_i = X_1 + X_2 +\cdots  +X_n$ where the covariances between the $X_i$'s are $\gamma_{ij}$,
the variance of $Y$ may be written as
%and their variances are $\gamma_{ii}$
%we have $ {\rm Var} (Y  ) =\sum_{i=1}^n  \gamma_{ii}^2 + C$ where $C \equiv \sum_{i \ne j}\gamma_{ij}$ is the contribution of correlations 
%to the variance of $Y$.  Moreover,
$\sigma_Y^2 \equiv  {\rm Var} (Y  ) = \sum_{i=1}^n \gamma_{i \cdot}$ where $\gamma_{i \cdot}=\sum_{j=1}^n \gamma_{ij}$
is the contribution of $X_i$ to the variance of $Y$.
The contribution to risk of $X_i$ to $Y$ is ${ \gamma_{i \cdot} / \sigma_Y }.$
\item Implied returns of a benchmark $b$: $ \lambda \Sigma b$.
\xei
\end{document}  


     \item What is the contribution of correlations to the risk of the mutual fund.\\
          {\em (a) 0\% (b) 20\% (c) 50\% (d) 70\% (e) 100\%}
          
          
\item The inverse of a $2 \times 2$ matrix:
$$
\left(\begin{array}{cc}
a & b\\
c & d 
\end{array}
\right)^{-1}
={1 \over ad-bc}
\left(\begin{array}{cc}
d & -b\\
-c & a 
\end{array}
\right)
$$


The first Quiz will be on Wednesday, 02/15/17. It will take about 30 minutes at the end of the class. The questions will be only on SAA topics covered through 02/01/17. All of the questions will be multiple choice. Some of them will be on concepts that we discussed in class and some will require short calculations � similar to what you had in HW1. Please review solutions to HW1 - they will be posted shortly after the deadline for the HW1 submission. All calculations, that you will need to do, will be short and can be done without a computer or calculator. You will not be allowed to use any electronic devices or books/notes. All the formulas, that you may need to answer the questions, will be provided.


Write	your	name	and	ID	on	every	sheet.	For	multiple-choice	questions,	unless	
otherwise	stated,	please	circle	the	correct	answer,	or,	if	necessary	for	clarity,	write	
your	choice	clearly	in	the	margin.	
This	quiz	is	worth	10%	of	your	total	grade	for	the	course.	You	have	30	minutes	to	
complete	the	quiz.	Finish	by	7:20 pm and	return	completed	quiz	paper	to	the	TAs